\section{平均値の定理}\label{節:平均値定理}
微分係数は1点での変化率である.一方,関数が区間全体で増えるかどうかを知るには,離れた2点の値を比べる必要がある.では,1点での傾きから,2点での値の差について何が言えるだろうか.

まずは簡単な関数で実験してみよう.区間の両端を結ぶ直線(割線)の傾きを求め,接線の傾き $f'(c)$ がそれに等しくなる点 $c$ を区間の内側で探すと,次のようになる.
\begin{center}
\begin{tabular}{cccc}
\toprule
$f(x)$ & 区間 & 割線の傾き & $c$\\
\midrule
$x^2$ & $[0,2]$ & $2$ & $1$\\
$x^3$ & $[0,3]$ & $9$ & $\sqrt3\approx1.73$\\
$x^3-x$ & $[-2,2]$ & $3$ & $\pm\frac{2}{\sqrt3}\approx\pm1.15$\\
\bottomrule
\end{tabular}
\end{center}
たとえば2行目では,割線の傾きは $\frac{27-0}{3-0}=9$ であり,\,$3c^2=9$ を解いて $c=\sqrt3$ を得る.図\ref{図:平均値定理}は1行目の場合を描いたものである.

\begin{figure}[htbp]
\centering
\begin{tikzpicture}[x=2.6cm,y=1.1cm,>=Stealth]
 \draw[->] (-0.15,0)--(2.35,0) node[right]{$x$};
 \draw[->] (0,-0.4)--(0,4.7) node[above]{$y$};
 \draw[very thick,domain=0:2.12,samples=70] plot (\x,{\x*\x});
 \draw[blue,thick] (0,0)--(2,4);
 \draw[red,thick,domain=0.4:1.8] plot (\x,{2*\x-1});
 \fill (0,0) circle(1.6pt);
 \fill (2,4) circle(1.6pt) node[above left]{$(2,4)$};
 \fill (1,1) circle(1.6pt) node[below right]{$(1,1)$};
 \draw[dashed] (1,0) node[below]{$1$}--(1,1);
 \draw[dashed] (2,0) node[below]{$2$}--(2,4);
 \node[blue] at (0.63,2.05) {割線};
 \node[red] at (1.8,2.2) {接線};
 \node at (1.35,4.55) {$y=x^2$};
\end{tikzpicture}
\caption{平均変化率に等しい微分係数をもつ点が,区間の内側に存在する.}
\label{図:平均値定理}
\end{figure}

どの例でも,割線と平行な接線が区間の内側で引けている.点 $c$ は区間の中央とは限らず,3行目のように1つに決まるとも限らない.そこで,「割線と傾きが等しい接線が,区間の内側のどこかで引ける」と予想できる.これを証明するために,まず極値で何が起こるかを確かめよう.

点 $c$ の十分近くのすべての $x$ で $f(x)\le f(c)$ なら,\,$f$ は $c$ で\textbf{極大値}をとるという.不等号を逆にしたものが\textbf{極小値}である.極大・極小は近くの点との比較であり,区間全体での最大・最小とは区別する.

$c$ が区間の内点で,\,$f$ がそこで微分可能かつ極大であるとしよう.\,$|h|$ が小さいとき $f(c+h)-f(c)\le0$ だから,差商の符号は $h$ の符号で決まる.
\[
 \frac{f(c+h)-f(c)}{h}\le0\quad(h>0),\qquad
 \frac{f(c+h)-f(c)}{h}\ge0\quad(h<0).
\]
$f$ は $c$ で微分可能なので,\,$h\to0+$ と $h\to0-$ の極限はともに $f'(c)$ に等しい.したがって $f'(c)\le0$ かつ $f'(c)\ge0$ となり,\,$f'(c)=0$ なのだ.極小の場合も不等号が逆になるだけである.ただし,端点では片側からしか比較できないので,この結論は使えない.

予想のうち,割線が水平な場合,つまり $f(a)=f(b)$ の場合から考えよう.両端の高さが等しければ,グラフは途中で山か谷を作るだろう.その頂上や谷底では,先ほどの議論から接線が水平になるはずである.

\begin{theorem}{ロルの定理}\label{定理:ロル}
$a<b$ とする.\,$f$ が閉区間 $[a,b]$ で連続,開区間 $(a,b)$ で微分可能で,\,$f(a)=f(b)$ なら,
\[
 f'(c)=0
\]
を満たす $c\in(a,b)$ が少なくとも1つ存在する.
\end{theorem}

定理\ref{定理:連続関数存在}により,\,$f$ は $[a,b]$ で最大値と最小値をとる.\,$f$ が定数なら,どの内点でも $f'=0$ である.定数でなければ,最大値と最小値の少なくとも一方は端点の値 $f(a)=f(b)$ と異なる.その値をとる点 $c$ は内点であり,\,$f$ はそこで極値をとるから,\,$f'(c)=0$ となる.

割線が水平でない場合は,割線の傾きの分だけグラフを傾け直してやれば,水平な場合に帰着する.

\begin{theorem}{平均値の定理}\label{定理:平均値}
$a<b$ とする.\,$f$ が閉区間 $[a,b]$ で連続,開区間 $(a,b)$ で微分可能なら,
\[
 f'(c)=\frac{f(b)-f(a)}{b-a}
\]
を満たす $c\in(a,b)$ が少なくとも1つ存在する.
\end{theorem}

割線の傾きを $m=\frac{f(b)-f(a)}{b-a}$ と書き,
\[
 g(x)=f(x)-m(x-a)
\]
と置く.\,$g$ は $f$ から割線と同じ傾きの一次関数を引いたものであり,
\[
 g(a)=f(a),\qquad g(b)=f(b)-\{f(b)-f(a)\}=f(a)
\]
となって両端の値が等しい.\,$g$ は $[a,b]$ で連続,\,$(a,b)$ で微分可能だから,ロルの定理により,ある $c\in(a,b)$ で
\[
 g'(c)=f'(c)-m=0
\]
となる.よって $f'(c)=m$ であり,予想は正しかった.

条件の役割も確認しておこう.\,$f(x)=|x|$ を $[-1,1]$ で考えると,両端の値は等しいが,微分可能な点での導関数は $-1$ か $1$ であり,\,$0$ にはならない.\,$x=0$ で微分できないため,ロルの定理も平均値の定理も適用できないのである.導関数の連続性は定理の仮定に含まれていないが,区間の内側のすべての点で微分可能であるという仮定は外せない.

\subsection{導関数と増減}
平均値の定理を使えば,導関数の符号から関数の増減が分かる.まず言葉を決めておこう.区間内の任意の2点 $x_1<x_2$ についてつねに $f(x_1)<f(x_2)$ となるとき,\,$f$ はその区間で\textbf{狭義に増加する}(狭義単調増加)という.つねに $f(x_1)>f(x_2)$ となるときは\textbf{狭義に減少する}という.そのどちらかであることを\textbf{狭義単調}という.\label{用語:狭義単調}「狭義」は,等号 $f(x_1)=f(x_2)$ を許さないという意味である.

区間内の任意の2点 $x_1<x_2$ に平均値の定理を使うと,その間のある $c$ に対して
\[
 f(x_2)-f(x_1)=f'(c)(x_2-x_1)
\]
となる.\,$x_2-x_1>0$ だから,区間内で常に $f'>0$ なら $f(x_2)>f(x_1)$ であり,\,$f$ は狭義に増加する.\,$f'<0$ なら狭義に減少する.\,$f'\ge0$ の場合は $f(x_2)\ge f(x_1)$ までしか言えず,値が一定の部分もあり得る.

同じ議論から,次の結果も得られる.
\begin{theorem}{導関数が同じ関数の差}\label{定理:導関数同じ}
1つの区間のすべての内点で $F'=G'$ が成り立ち,端点を含む場合にはそこで両関数が連続なら,その区間で $F-G$ は定数である.特に,導関数が恒等的に $0$ なら,元の関数は定数である.
\end{theorem}

実際,\,$H=F-G$ と置くと $H'=0$ であり,上の式から任意の2点について $H(x_2)-H(x_1)=0$ となる.後の積分では,微分して同じ関数になる候補が,定数の違いを除いて同じであることをこの定理で確かめる.

ここで「1つの区間」という条件は外せない.たとえば
\[
 H(x)=
 \begin{cases}
  0 & (x<0)\\
  1 & (x>0)
 \end{cases}
\]
と定めれば,定義域のすべての点で $H'=0$ だが,\,$H$ は定数ではない.定義域が $x<0$ と $x>0$ の2つの区間に分かれているので,それぞれの区間で別の定数をとれるのである.
%
